\documentclass[11pt]{article}
\usepackage[margin=1in]{geometry}
\usepackage{amsmath,amssymb}
\usepackage{booktabs}
\usepackage{hyperref}
\usepackage{xcolor}

\title{Replication Report: Bridi \& Marquezino (2024)\\
\large ``Analytical results for the Quantum Alternating Operator Ansatz with Grover Mixer''}
\author{Ollie (OpenClaw AI) --- QC-100 Replication Project}
\date{2026-07-03 (backfilled 2026-07-06)}

\begin{document}
\maketitle

\begin{abstract}
We reproduce the checkable content of Bridi \& Marquezino, arXiv:2401.11056v3 (12 Aug 2024). The paper's central analytical premise --- that Grover-mixer QAOA (GM-QAOA) is \emph{permutation-invariant} in the cost-function domain, and therefore ``structurally blind'' --- is verified to $10^{-15}$ across 20 random permutations, two graphs, and two depths. Equation (8) of the paper (closed-form marked-state probability for a binary cost) is reproduced to machine precision for $r \in \{1,2,3,4\}$ on $n=6$ qubits. On a structured 6-node MAX-CUT the standard transverse-field (X-mixer) QAOA beats GM-QAOA at $p=1,2,3$ with a growing approximation-ratio gap ($+0.068, +0.141, +0.147$), reproducing the paper's Introduction claim that structural blindness is a real performance cost. The paper's main asymptotic theorem (exponential-in-$n$ round count on complete-bipartite MAX-CUT) is \emph{not} tested here --- it is inherently a scaling result.

Verdict: \textbf{REPLICATED} for all four small-instance / closed-form claims exercised.
\end{abstract}

\section{Paper and headline claim}
Bridi \& Marquezino formalize GM-QAOA, whose mixer
\[
U_{\mathrm{GM}}(\beta) \;=\; \exp\!\bigl(-\,\mathrm{i}\,\beta\,|s\rangle\!\langle s|\bigr)
\;=\; I + \bigl(e^{-\mathrm{i}\beta}-1\bigr)|s\rangle\!\langle s|,
\qquad |s\rangle = H^{\otimes n}|0\rangle,
\]
is a rank-1 reflection about the uniform superposition. Their central structural observation is that $\langle C\rangle_{\mathrm{GM}}$ depends only on the \emph{multiset} of cost values, not on which bitstring holds which value --- so any GM-based QAOA can at best match Grover-style quadratic search. As a corollary, on complete-bipartite MAX-CUT the number of rounds must grow exponentially in vertex count to guarantee any fixed approximation ratio.

\section{Headline claim exercised in this replication}
The most \emph{directly checkable} claims on public artifacts alone (statevector, small $n$):
\begin{itemize}
  \item \textbf{C1 (headline structural claim, exercised)}: GM-QAOA $\langle C\rangle$ is invariant under any permutation of the cost domain.
  \item \textbf{C2 (contrast, exercised)}: the same identity must \emph{fail} for the transverse-field X-mixer.
  \item \textbf{C3 (Eq.~(8), exercised)}: for binary cost with marked-fraction $\rho \le \rho_{\mathrm{Th}}(r)$,
  \[
    P(\rho,r) = \sin^2\!\bigl((2r+1)\arcsin\sqrt{\rho}\bigr),\qquad \rho_{\mathrm{Th}}(r) = \sin^2\!\tfrac{\pi}{4r+2}.
  \]
  \item \textbf{C4 (introductory advantage claim, exercised)}: X-mixer beats GM-QAOA at fixed depth on structured MAX-CUT.
  \item \textbf{C5 (main asymptotic theorem, NOT exercised)}: exponential-in-$n$ round requirement on complete-bipartite MAX-CUT.
\end{itemize}

\section{Method}
All experiments are pure-Python statevector simulations at $n=5,6$ qubits (state dim $32,64$).
Environment: Python 3.14.6, \texttt{numpy 2.5.0}, \texttt{scipy 1.18.0}, \texttt{qiskit 2.5.0} (Qiskit installed for cross-check, but the reported numbers come from the NumPy simulator so the physics is fully under our control). Free endpoints only; no external LLM in-the-loop for numerical claims.

\section{Results}

\subsection{C1, C2 --- permutation (in)variance}
\begin{center}
\begin{tabular}{lccrr}
\toprule
Graph & $p$ & Mixer & $\max_\pi\lvert\langle C\rangle_\pi - \langle C\rangle_{\mathrm{id}}\rvert$ & std over perms \\
\midrule
A ($n=6$, 8 edges) & 2 & GM & $8.88\times 10^{-16}$ & $1.44\times 10^{-15}$ \\
A ($n=6$, 8 edges) & 2 & X  & $1.33$                 & $3.11\times 10^{-1}$ \\
B ($n=5$, 7 edges) & 3 & GM & $1.78\times 10^{-15}$ & $2.63\times 10^{-15}$ \\
B ($n=5$, 7 edges) & 3 & X  & $2.57$                 & $5.9\times 10^{-1}$ \\
\bottomrule
\end{tabular}
\end{center}
GM-QAOA is permutation-invariant to floating-point precision across every one of 20 random permutations, at two depths and two graphs. The X-mixer is manifestly non-invariant.

\subsection{C3 --- Grover-binary closed form (Eq.~(8))}
\begin{center}
\begin{tabular}{rrrrrrr}
\toprule
$r$ & $k$ & $\rho=k/64$ & $\rho_{\mathrm{Th}}(r)$ & $P_{\mathrm{stat}}$ & $P_{\mathrm{Eq.(8)}}$ & $\lvert\Delta\rvert$ \\
\midrule
1 & 15 & 0.234375 & 0.250000 & 0.9970095 & 0.9970095 & $1.11\times 10^{-16}$ \\
2 &  6 & 0.093750 & 0.095492 & 0.9997793 & 0.9997793 & 0 \\
3 &  3 & 0.046875 & 0.049516 & 0.9981394 & 0.9981394 & $1.11\times 10^{-16}$ \\
4 &  1 & 0.015625 & 0.030154 & 0.8163772 & 0.8163772 & 0 \\
\bottomrule
\end{tabular}
\end{center}
Machine precision agreement at every depth tested.

\subsection{C4 --- X-mixer vs GM-QAOA on structured MAX-CUT}
Graph A, $\text{cost}_{\max}=8$, uniform-random baseline ratio $0.500$.
\begin{center}
\begin{tabular}{crrrrr}
\toprule
$p$ & GM $\langle C\rangle$ & GM ratio & X $\langle C\rangle$ & X ratio & X$-$GM ratio \\
\midrule
1 & 5.086 & 0.6357 & 5.633 & 0.7041 & $+0.068$ \\
2 & 5.899 & 0.7374 & 7.030 & 0.8787 & $+0.141$ \\
3 & 6.496 & 0.8120 & 7.673 & 0.9591 & $+0.147$ \\
\bottomrule
\end{tabular}
\end{center}
Gap grows with $p$, consistent with the paper's structural-blindness argument.

\section{Critique --- what we \emph{did} and did \emph{not} exercise}
\textbf{What is genuine here.} We independently reimplemented GM-QAOA \emph{from its rank-1 exact form}, not by importing a black-box GM-QAOA package: $U_{\mathrm{GM}}(\beta) = I + (e^{-\mathrm{i}\beta}-1)|s\rangle\!\langle s|$. The permutation-invariance identity is a strict linear-algebra prediction of that operator, and it comes out to $10^{-15}$ --- floating-point noise --- across 20 random permutations. That is a strong reimplementation-level replication of the paper's central analytical premise, and it is what the paper's whole ``structural blindness'' story hangs on.
The Eq.~(8) test is likewise a first-principles closed-form check ($P = \sin^2((2r+1)\arcsin\sqrt{\rho})$) against a statevector we constructed ourselves; agreement to $10^{-16}$ means the Grover-mixer physics in our code is the same physics as in the paper's derivation.

\textbf{Where the replication is thinner than the paper.} Three honest caveats:
\begin{enumerate}
  \item \emph{The paper's main new theorem (C5) is asymptotic and is not tested.} A single small-instance CPU simulation cannot falsify or confirm ``exponential in $n$'' behavior on complete-bipartite MAX-CUT --- that requires a scaling sweep across many $n$ and a log-slope fit. We flag this explicitly rather than paper over it.
  \item \emph{The X-vs-GM comparison in C4 is on \emph{one} 6-node graph.} The paper's Introduction cites this as a general phenomenon on structured problems. We saw a clean growing gap on Graph A, but we did not sweep graph families (regular graphs, Erd\H{o}s--R\'enyi, bipartite, weighted MAX-CUT). It is possible the gap sign is graph-dependent, or that with more COBYLA restarts GM-QAOA closes some of it. Our result is consistent with the paper's claim but is one instance, not a distribution.
  \item \emph{The threshold variant GM-Th-QAOA (Theorem~1) is not implemented.} Its closed form would be another clean check, but it needs the threshold phase separator, which we did not build. The core GM-mixer physics (C1, C2, C3) is exercised, so this is a coverage gap rather than a physics gap.
\end{enumerate}

\textbf{Reimplemented independently vs. quoted.} We did \emph{not} take the paper's numbers on faith. For C1--C4 we produced our own numbers from our own simulator built from the operator definition, and the paper's central quantitative predictions (permutation invariance, Eq.~(8), and the direction/growth of the X-vs-GM gap) were reproduced.

\textbf{Comparison against standard QAOA baseline.} Yes, explicitly: C2 (X-mixer non-invariance) and C4 (X-mixer ratio at $p=1,2,3$) are the required X-mixer baseline. The mixer-choice advantage held quantitatively for the X-mixer on our structured instance, matching the paper's introductory statement and its refs~[13,14]. The paper's own \emph{claim} is a mixer-choice \emph{disadvantage} for GM-QAOA on structured problems; that is what we saw.

\section{Verdict}
\textbf{REPLICATED} for C1, C2, C3, C4 (headline structural claim + Eq.~(8) closed form + X-vs-GM comparison on a structured graph). C5 (asymptotic theorem) is out of scope for a single-instance replication and is flagged as untested rather than claimed. See \texttt{report/failure\_analysis.md} for a fuller critique and \texttt{report/open\_questions.json} for follow-ups.

\end{document}
